\documentclass[11pt,reqno]{amsart}
\usepackage[T1]{fontenc}
\usepackage[utf8]{inputenc}
\usepackage{lmodern}
\usepackage[margin=1in]{geometry}
\usepackage{amsmath,amssymb,amsthm,mathtools}
\usepackage{booktabs,array,longtable}
\usepackage{microtype}
\usepackage{xcolor}
\usepackage{enumitem}
\usepackage{listings}
\usepackage{xurl}
\usepackage[colorlinks=true,linkcolor=blue!45!black,citecolor=blue!45!black,urlcolor=blue!45!black]{hyperref}
\hypersetup{pdftitle={Posterior cones and sharp reconstruction thresholds on finite-mean Galton--Watson trees},pdfauthor={Yitzchak Shmalo}}
\setlength{\parskip}{0.25em}
\setlength{\parindent}{1.2em}
\emergencystretch=3em
\setlist[itemize]{topsep=4pt,itemsep=2pt}
\setlist[enumerate]{topsep=4pt,itemsep=2pt}
\allowdisplaybreaks[2]
\newtheorem{theorem}{Theorem}[section]
\newtheorem{proposition}[theorem]{Proposition}
\newtheorem{lemma}[theorem]{Lemma}
\newtheorem{corollary}[theorem]{Corollary}
\theoremstyle{definition}
\newtheorem{definition}[theorem]{Definition}
\theoremstyle{remark}
\newtheorem{remark}[theorem]{Remark}
\newcommand{\E}{\mathbb E}
\newcommand{\Prob}{\mathbb P}
\newcommand{\Pbb}{\mathbb P}
\newcommand{\TV}{\operatorname{TV}}
\newcommand{\eps}{\varepsilon}
\newcommand{\R}{\mathbb R}
\newcommand{\N}{\mathbb N}
\newcommand{\ind}{\mathbf 1}
\newcommand{\cC}{\mathcal C}
\newcommand{\KL}{\operatorname{KL}}
\newcommand{\supp}{\operatorname{supp}}
\newcommand{\Law}{\mathcal L}
\lstset{basicstyle=\ttfamily\footnotesize,breaklines=true,columns=fullflexible,keepspaces=true,frame=single,rulecolor=\color{black!20},showstringspaces=false,tabsize=2}
\title[Posterior cones and reconstruction]{Posterior cones and sharp reconstruction thresholds on finite-mean Galton--Watson trees}
\author{Yitzchak Shmalo}
\address{Einstein Institute of Mathematics, The Hebrew University of Jerusalem, Givat Ram, Jerusalem, Israel}
\email{yitzchak.shmalo@gmail.com}
\date{24 October 2026}

\begin{document}
\begin{abstract}
We study reconstruction for symmetric ferromagnetic Potts broadcasting on an observed Galton--Watson tree with arbitrary finite mean offspring. For three states we construct an explicit invariant condition on posterior laws using the polynomial $H=y-x^2/2$, where $x$ and $y$ are the quadratic and cubic centered posterior invariants. Two global inequalities, certified by exact nonnegative coefficients, preserve this condition and give strict information contraction. Combining the resulting argument with an existing four-state posterior cone proves that reconstruction occurs exactly when $d\lambda^2>1$ for both three and four states, including nonreconstruction at equality. No offspring variance or tail assumption is required. We give explicit critical decay bounds, a first-moment-only supercritical argument, and almost-sure statements conditional on survival. Exact verification programs and complete coefficient tables accompany the proof. The noiseless unary-line endpoint is classified separately.
\end{abstract}
\keywords{Reconstruction on trees, Potts model, Kesten--Stigum threshold, Galton--Watson process, posterior experiments, exact polynomial certificates}
\maketitle

\section{Introduction and main result}

The reconstruction problem asks whether a noisy broadcast retains information about its root at arbitrarily large distances. For the symmetric $q$-state channel
\begin{equation}\label{eq:channel}
 K_\lambda(i,j)=\lambda\ind_{\{i=j\}}+\frac{1-\lambda}{q},
 \qquad i,j\in\{1,\ldots,q\},
\end{equation}
the Kesten--Stigum parameter on a tree of mean offspring $d$ is $d\lambda^2$. Reconstruction above this parameter's unit threshold follows from a second-moment argument. The difficult direction is to rule out nonlinear reconstruction at and below the threshold.

Sly proved sharpness for the three-state model at sufficiently large regular degree and non-sharpness for every regular degree when $q\ge5$ \cite{Sly2011}. More recent manuscripts establish the three-state symmetric and four-state ferromagnetic thresholds on all regular and observed Poisson trees \cite{OpenAI3,OpenAI4}. The positive-channel argument in the three-state manuscript uses offspring-specific entropy estimates, whereas the four-state manuscript supplies a cone of posterior laws stable under pairwise products.

This paper gives a simple three-state cone with global exact certificates and develops the finite-mean Galton--Watson consequences for both three and four states. The three-state certificate is self-contained below and in Appendix~\ref{app:exact-certificates}. The four-state invariant and its pair inequalities are those of \cite{OpenAI4}. Combining these two cones yields a unified proof for arbitrary finite-mean offspring distributions, with explicit decay estimates at the critical threshold.

Let $D$ be a nonnegative integer-valued random variable with finite mean $d$, and let $T$ be a Galton--Watson tree with iid offspring $D$. The tree is independent of the spin process. The root spin $\sigma_\rho$ is uniform on $\{1,\ldots,q\}$; conditional on a parent spin, its children are generated independently through \eqref{eq:channel}. The observer knows $T$ and the spins on level $L_n(T)$. Put
\begin{equation}\label{eq:advantage}
 p_n=\Prob(\sigma_\rho=\cdot\mid T,\sigma_{L_n(T)}),\qquad
 a_n(T)=\E_\sigma[\TV(p_n,u_q)\mid T],\qquad
 a_n=\E_T a_n(T),
\end{equation}
where $u_q=(1/q,\ldots,1/q)$. Reconstruction means that the corresponding limiting advantage is positive. For a fixed locally finite tree, $a_n(T)$ is nonincreasing by data processing, so this limit exists.

\begin{theorem}[Finite-mean reconstruction threshold]\label{thm:main}
Let $q\in\{3,4\}$ and $0<\lambda<1$. For symmetric Potts broadcasting on the observed Galton--Watson tree described above:
\begin{enumerate}
\item If $d\lambda^2\le1$, then $a_n(T)\to0$ for almost every tree.
\item If $d\lambda^2>1$, then $\lim_n a_n(T)>0$ for almost every surviving tree, and the limit is zero on extinction. In particular $\lim_n a_n>0$.
\end{enumerate}
Thus annealed reconstruction occurs if and only if $d\lambda^2>1$. The result requires only $\E D<\infty$ and includes every regular $d$-ary tree with integer $d\ge2$.
\end{theorem}

The exclusion of $\lambda=1$ from the finite-noise theorem is substantive only at the deterministic unary tree: $D\equiv1,\lambda=1$ reconstructs perfectly although $d\lambda^2=1$. At $\lambda=1$ and $d>1$ reconstruction is perfect on survival; nondegenerate critical and subcritical Galton--Watson trees become extinct almost surely. The case $\lambda=0$ is immediate. All results here concern the ferromagnetic regime; no antiferromagnetic conclusion is obtained by the cone argument.

\section{Posterior experiments and Bayes products}\label{sec:experiments}

An experiment with uniform input $X\in\{1,\ldots,q\}$ induces a marginal law $\mu$ of the posterior vector $p=\Prob(X=\cdot\mid Y)$ in the simplex $\Delta_q$. Its barycenter is $u_q$. We call the law symmetric if it is invariant under every permutation of the input labels.

\begin{lemma}[Bayes product]\label{lem:bayes}
Suppose two experiments are independent conditional on the common uniform input, and their marginal posterior laws are $\mu$ and $\nu$. For independent $p\sim\mu$ and $r\sim\nu$, define
\[
 U_i=q p_i r_i,\qquad z=\sum_iU_i,\qquad P_i=U_i/z\quad(z>0).
\]
For every bounded measurable $f$,
\begin{equation}\label{eq:bayes-product}
 \E_{\mu\star\nu}f(P)=\E_{\mu\otimes\nu}[z f(P)],
\end{equation}
with zero contribution when $z=0$. If the laws are symmetric and $f$ is permutation invariant, the integrand may be averaged over every relative permutation of $r$.
\end{lemma}

\begin{proof}
With respect to each observation's marginal law, its density conditional on $X=i$ is $q p_i$ or $q r_i$. Conditional independence gives joint marginal density $\frac1q\sum_i(qp_i)(qr_i)=z$ relative to the product of the two marginals. Bayes' formula gives the displayed posterior and \eqref{eq:bayes-product}. The marginal posterior barycenters imply $\E z=1$. If a normalizer is zero, that pair has zero probability under the joint observation law. Symmetry permits averaging the second marginal over relative permutations without changing its distribution.
\end{proof}

It is essential that independence in \eqref{eq:bayes-product} is under the reference product measure. The actual joint observation measure contains the factor $z$.

\begin{lemma}[Edges, mixtures, and additional observations]\label{lem:operations}
The following operations on posterior experiments have the stated properties.
\begin{enumerate}
\item Passing a child experiment through one channel edge maps its posterior to
\begin{equation}\label{eq:edge}
 p\longmapsto u_q+\lambda(p-u_q).
\end{equation}
\item An observed mixing label independent of the input mixes the corresponding posterior laws.
\item If $\Phi$ is convex on $\Delta_q$, observing additional information cannot decrease $\E\Phi(p)$.
\end{enumerate}
\end{lemma}

\begin{proof}
The first assertion follows by applying the symmetric transition matrix to the child posterior; the uniform child marginal equals the uniform parent marginal. The second is conditioning on the observed label. For the third, the posterior before the added observation is the conditional expectation of the posterior afterward. Conditional Jensen gives the claim.
\end{proof}

\section{An explicit three-state posterior cone}\label{sec:q3}

For $p\in\Delta_3$, let $v_i=3p_i-1$ and define
\begin{equation}\label{eq:q3-invariants}
 x(p)=\frac16\sum_i v_i^2=\frac{3\sum_i p_i^2-1}{2},
 \qquad y(p)=\frac16\sum_i v_i^3.
\end{equation}
The functions used throughout this section are
\begin{equation}\label{eq:q3-potential}
 \boxed{H(p)=y(p)-\tfrac12x(p)^2,\qquad
 F(p)=\tfrac32x(p)-y(p),\qquad G(p)=1-x(p).}
\end{equation}

\begin{lemma}[Scalar properties]\label{lem:q3-scalar}
For all $p\in\Delta_3$,
\[
 0\le x(p)\le1,\qquad |y(p)|\le x(p)^{3/2}\le x(p),
 \qquad F(p)\ge\tfrac12x(p)\ge0,\qquad G(p)\ge0.
\]
The revealing posterior has $H=1/2$, and the uniform posterior has $H=0$. Under \eqref{eq:edge},
\begin{equation}\label{eq:q3-radial}
 H\bigl(u_3+\lambda(p-u_3)\bigr)
 =\lambda^3H(p)+\tfrac12\lambda^3(1-\lambda)x(p)^2
 \ge\lambda^3H(p)\quad(0\le\lambda\le1).
\end{equation}
\end{lemma}

\begin{proof}
Set $\omega=e^{2\pi i/3}$ and $\zeta=p_1+\omega p_2+\omega^2p_3$. Direct expansion gives $x=|\zeta|^2$ and $y=\operatorname{Re}(\zeta^3)$. Since $|\zeta|\le1$, all bounds follow. At a point mass $x=y=1$, whereas both vanish at $u_3$. Equation~\eqref{eq:edge} multiplies $\zeta$ by $\lambda$, so $x$ and $y$ scale by $\lambda^2$ and $\lambda^3$. Substitution yields \eqref{eq:q3-radial}.
\end{proof}

For $p,r\in\Delta_3$ and $\pi\in S_3$, write
\begin{equation}\label{eq:q3-orbits}
 z_\pi=3\sum_i p_i r_{\pi(i)},\qquad
 P_i^\pi=\frac{3p_i r_{\pi(i)}}{z_\pi}\quad(z_\pi>0),
 \qquad [zf(P)]=\frac16\sum_{\pi\in S_3}z_\pi f(P^\pi).
\end{equation}
Terms with $z_\pi=0$ are zero.

\begin{proposition}[Global three-state pair inequalities]\label{prop:q3-pair}
With $\eps=1/1000$, every $p,r\in\Delta_3$ satisfy
\begin{align}
 [zx(P)]&\le x(p)+x(r)-H(p)F(r)-F(p)H(r)-\eps x(p)x(r),\label{eq:q3-pair-A}\\
 [zH(P)]&\ge H(p)G(r)+G(p)H(r)+\eps x(p)x(r).\label{eq:q3-pair-B}
\end{align}
\end{proposition}

\begin{proof}
We give the exact polynomial certificate and its domain argument. The independent integer implementation and complete verification data are described in Appendix~\ref{app:exact-certificates}.

The expressions are invariant under separate permutations of $p$ and $r$. A sorted vector $p_1\ge p_2\ge p_3$ is a convex combination of
\[
 e=(1,0,0),\qquad h=(1/2,1/2,0),\qquad u=(1/3,1/3,1/3),
\]
with coefficients $p_1-p_2$, $2(p_2-p_3)$, and $3p_3$. Thus it suffices to work on the product of these two ordered chambers, in nonnegative barycentric triples each summing to one.

For a fixed permutation define
\[
 \begin{gathered}
 U_i=3p_i r_{\pi(i)},\qquad Z_\pi=\sum_iU_i,\qquad V_i=3U_i-Z_\pi,\\
 X_\pi=\tfrac16\sum_iV_i^2,\qquad Y_\pi=\tfrac16\sum_iV_i^3,
 \qquad N_\pi=Y_\pi Z_\pi-\tfrac12X_\pi^2.
 \end{gathered}
\]
Direct substitution gives $Z_\pi x(P^\pi)=X_\pi/Z_\pi$ and
$Z_\pi H(P^\pi)=N_\pi/Z_\pi^3$. Let
\begin{align*}
 K_A&=x(p)+x(r)-H(p)F(r)-F(p)H(r)-\eps x(p)x(r),\\
 K_B&=-H(p)G(r)-G(p)H(r)-\eps x(p)x(r).
\end{align*}
The denominator-cleared targets are
\begin{align}
 T_A&=6K_A\prod_\pi Z_\pi-\sum_\pi X_\pi\prod_{\tau\ne\pi}Z_\tau,\label{eq:q3-cleared-A}\\
 T_B&=6K_B\prod_\pi Z_\pi^3+\sum_\pi N_\pi\prod_{\tau\ne\pi}Z_\tau^3.\label{eq:q3-cleared-B}
\end{align}
Homogenize each coordinate block to degree $10$ for $T_A$ and degree $22$ for $T_B$. Multiplication by powers of the coordinate sums leaves values unchanged on the simplex. The supplied exact integer programs produce the following exhaustive coefficient counts:
\begin{center}
\begin{tabular}{lrrrr}
\toprule
Target & Total & Positive & Zero & Negative\\
\midrule
$T_A$ & 4,356 & 3,873 & 483 & 0\\
$T_B$ & 76,176 & 71,627 & 4,549 & 0\\
\bottomrule
\end{tabular}
\end{center}
Every coefficient is rational with positive denominator. Nonnegative coefficients prove $T_A,T_B\ge0$ for all nonnegative barycentric coordinates. This is a global certificate on the initial chamber; no subdivision or sampled extrapolation is required. A separately written sparse six-variable implementation regenerates every coefficient and agrees exactly after the explicitly documented positive integer scaling.

For strictly positive $p,r$, all $Z_\pi$ are positive, so division recovers \eqref{eq:q3-pair-A}--\eqref{eq:q3-pair-B}. On the boundary, $x$ and $H$ are bounded on the compact simplex. Hence $Z_\pi f(P^\pi)\to0$ whenever $Z_\pi\to0$. Approximating each boundary pair by an interior pair proves the stated inequalities everywhere.
\end{proof}

\begin{definition}
Let $\cC_3$ be the symmetric posterior laws $\mu$ such that $\E_\mu H\ge0$.
\end{definition}

The condition is imposed on an expectation, not on the support of the law. For example, $H(1/2,1/2,0)=-5/32$; such a posterior value is not excluded from an experiment whose average $H$ is nonnegative.

\begin{proposition}[Closure and strict subadditivity]\label{prop:q3-cone}
The class $\cC_3$ contains both full revelation and the uninformative experiment. It is preserved by positive channel steps, observed input-independent mixtures, and conditionally independent products. If $\mu,\nu\in\cC_3$ and $m(\mu)=\E_\mu x$, then
\begin{equation}\label{eq:q3-subadditivity}
 m(\mu\star\nu)\le m(\mu)+m(\nu)-\eps m(\mu)m(\nu).
\end{equation}
\end{proposition}

\begin{proof}
The initial laws and edge closure follow from Lemma~\ref{lem:q3-scalar}. Mixture closure follows by linearity and symmetry. For products, integrate \eqref{eq:q3-pair-B} against $\mu\otimes\nu$ and use Lemma~\ref{lem:bayes}:
\[
 \E_{\mu\star\nu}H\ge
 (\E_\mu H)(\E_\nu G)+(\E_\mu G)(\E_\nu H)
 +\eps (\E_\mu x)(\E_\nu x)\ge0.
\]
Product experiments remain symmetric. Integrating \eqref{eq:q3-pair-A}, the two products involving $H$ and $F$ are nonnegative because $F\ge0$ and both $H$ expectations are nonnegative. Dropping them proves \eqref{eq:q3-subadditivity}.
\end{proof}

% Imported q=4 cone and the finite-mean Galton--Watson arguments.
% Uses only standard LaTeX/amsmath/amssymb/amsthm commands.
% Required theorem environments: proposition, lemma, theorem, remark.
% Bibliography key: OpenAI4.

\section{Four states: a posterior cone}
\label{sec:four-state-cone}

We use the four-state posterior cone and global pair inequalities of \cite{OpenAI4} and derive their consequences for arbitrary finite-mean offspring. We state the inequalities and describe their finite certificates explicitly.

For a posterior $p\in\Delta_4$, let
\[
 \zeta=Lp,\qquad
 L=\begin{pmatrix}
  1&1&-1&-1\\
  1&-1&1&-1\\
  1&-1&-1&1
 \end{pmatrix}.
\]
Define
\begin{align}
 A&=\zeta_1^2+\zeta_2^2+\zeta_3^2
      =4\sum_{i=1}^4p_i^2-1,
 & B&=\zeta_1\zeta_2\zeta_3,
 & C&=\zeta_1^2\zeta_2^2+\zeta_1^2\zeta_3^2+
       \zeta_2^2\zeta_3^2.
 \label{eq:q4-basic-polynomials}
\end{align}
In particular $0\le A\le3$.  The three polynomials needed below are
\begin{align}
 H_4={}&B-\frac1{10}A^2-\frac3{20}C+\frac3{10}AB
       -\frac1{50}A^3+\frac7{100}AC-\frac{51}{100}B^2,
       \label{eq:q4-H}\\
 F_4={}&4A-3B+\frac7{10}A^2-4C-\frac45AB,
       \label{eq:q4-F}\\
 G_4={}&1-\frac{161}{100}A-\frac{31}{10}B
       +\frac{773}{1000}A^2+4C-\frac{199}{50}AB.
       \label{eq:q4-G}
\end{align}
We regard these both as functions of $p$ and as polynomials in $\zeta=Lp$;
the argument will indicate which representation is used.

\subsection{The four-state inequalities}

For $p,r\in\Delta_4$ and $\pi\in S_4$, write
\[
 U_i^\pi=4p_i r_{\pi(i)},\qquad
 z_\pi=\sum_{i=1}^4U_i^\pi,\qquad
 P_i^\pi=\frac{U_i^\pi}{z_\pi}\quad(z_\pi>0).
\]
For a bounded function $f$ on $\Delta_4$, use the notation
\[
 [zf(P)]_4(p,r)=\frac1{24}\sum_{\pi\in S_4}z_\pi f(P^\pi),
\]
where a summand with $z_\pi=0$ is defined to be zero.  Thus the brackets
average over all twenty-four relative color permutations.

\begin{proposition}[Four-state polynomial inequalities, \cite{OpenAI4}]
\label{prop:q4-imported-pair}
The polynomials \eqref{eq:q4-H}--\eqref{eq:q4-G} satisfy
\[
 H_4\ge-\frac12,\qquad F_4\ge0,\qquad G_4\ge0
 \quad\hbox{on }\Delta_4.
\]
For $0\le t\le1$ and $\zeta\in L\Delta_4$,
\begin{equation}
 H_4(t\zeta)\ge t^3H_4(\zeta).
 \label{eq:q4-radial}
\end{equation}
For every $p,r\in\Delta_4$,
\begin{align}
 [zA(P)]_4(p,r)
 &\le A(p)+A(r)-H_4(p)F_4(r)-F_4(p)H_4(r)\notag\\
 &\hspace{22mm}
 -\frac{A(p)A(r)\bigl(A(p)+A(r)\bigr)}{1000},
 \label{eq:q4-pair-A}\\
 [zH_4(P)]_4(p,r)
 &\ge H_4(p)G_4(r)+G_4(p)H_4(r).
 \label{eq:q4-pair-H}
\end{align}
\end{proposition}

\subsection{Finite certificates for the four-state inequalities}
\label{subsec:q4-certificate}

We summarize the exact verification in \cite{OpenAI4} to specify precisely
the four-state computational ingredient.  Permuting the four probabilities
acts on $Lp$ by a signed permutation of three coordinates, with the product
of the three signs equal to one.  Consequently $A,B,C$ and
$H_4,F_4,G_4$ are invariant.  The averaged pair expressions are invariant
under separate permutations of $p$ and $r$.

It therefore suffices to work on the ordered simplex
\[
 \mathcal S_4=\{p\in\Delta_4:p_1\ge p_2\ge p_3\ge p_4\},
\]
whose vertices are
\[
 e_1,\quad \frac{e_1+e_2}{2},\quad
 \frac{e_1+e_2+e_3}{3},\quad u_4.
\]
Indeed its barycentric coordinates are
$p_1-p_2,2(p_2-p_3),3(p_3-p_4),4p_4$.
On a simplex, a homogeneous polynomial with nonnegative ordinary monomial
coefficients is nonnegative.  Homogenization by powers of the sum of the
barycentric coordinates does not alter values on the simplex.  This gives
the following exact scalar certificates; the displayed minimum refers to
ordinary homogeneous coefficients.
\begin{center}
\begin{tabular}{lrrr}
\hline
Polynomial & Degree & Coefficients & Minimum\\
\hline
$H_4+1/2$ & 6 & 84 & $19/50$\\
$F_4$ & 5 & 56 & $0$\\
$G_4$ & 14 & 680 & $87/1000$\\
$(3-\mathcal D)H_4$ & 6 & 84 & $0$\\
\hline
\end{tabular}
\end{center}
Here $\mathcal D=\sum_{j=1}^3\zeta_j\partial_{\zeta_j}$ is the Euler
operator.  Its scalar certificate implies \eqref{eq:q4-radial}, because
\[
 \frac{d}{dt}\bigl(t^{-3}H_4(t\zeta)\bigr)
 =t^{-4}(\mathcal D-3)H_4(t\zeta)\le0
 \qquad(0<t\le1),
\]
and $H_4(0)=0$ handles $t=0$.

For the pair inequalities, define
\begin{align*}
 D_0(p,r)&=\frac{A(p)A(r)(A(p)+A(r))}{1000},\\
 K_a&=24\bigl(A(p)+A(r)-H_4(p)F_4(r)-F_4(p)H_4(r)-D_0(p,r)\bigr),\\
 K_b&=-24\bigl(H_4(p)G_4(r)+G_4(p)H_4(r)\bigr).
\end{align*}
The two targets have the common form
\begin{equation}
 \mathcal R=K+\sum_{\pi\in S_4}\frac{N_\pi}{z_\pi^k},
 \label{eq:q4-rational-target}
\end{equation}
with the following data:
\begin{center}
\begin{tabular}{cclc}
\hline
Target & $k$ & $N_\pi$ & $m$\\
\hline
$a$ & 1 & $-\lVert LU^\pi\rVert_2^2$ & $3$\\
$b$ & 5 & $z_\pi^6H_4(U^\pi/z_\pi)$ & $1/2$\\
\hline
\end{tabular}
\end{center}
The numerator in the second row is polynomial: each degree-$s$ term of
$H_4$ is multiplied by $z_\pi^{6-s}$ after substitution.  The two values of
$\mathcal R$ are exactly twenty-four times the respective nonnegative
differences in \eqref{eq:q4-pair-A} and \eqref{eq:q4-pair-H}.

The verifier covers $\mathcal S_4\times\mathcal S_4$ by products of rational
subsimplices.  On each cell it applies one of two sufficient, exact
coefficient tests.

The first test uses
\[
 t_k(z)=\sum_{j=0}^5\binom{-k}{j}(z-1)^j\le z^{-k}
 \qquad(z>0).
\]
Taylor's theorem proves the inequality on the entire positive half-line:
the sixth derivative of $z^{-k}$ is positive, and the remainder contains
the nonnegative factor $(z-1)^6$.  The scalar bounds above give
$N_\pi+m z_\pi^{k+1}\ge0$.  Hence
\[
 \frac{N_\pi}{z_\pi^k}
 \ge (N_\pi+m z_\pi^{k+1})t_k(z_\pi)-m z_\pi.
\]
Nonnegative homogeneous coefficients of the sum of these polynomial lower
bounds and $K$ certify the cell.  The nonnegative multiplier is necessary;
one cannot multiply a lower bound by an uncontrolled signed numerator.

For the second test, choose rational bihomogeneous polynomials $J_\pi$ of
bidegree $(3,3)$.  With
$W_j(X,Y)=(\sum_iX_i)^j(\sum_iY_i)^j$, form
\[
 O_\pi=W_2N_\pi-z_\pi^kJ_\pi,
 \qquad Q_\pi=z_\pi^kW_3.
\]
These are both of bidegree $(k+3,k+3)$, and $Q_\pi$ has nonnegative
coefficients.  The checker chooses a rational $c_\pi\ge0$ such that every
coefficient of $O_\pi+c_\pi Q_\pi$ is nonnegative.  A negative coefficient
of $O_\pi$ at a zero coefficient of $Q_\pi$ makes this test fail.  Otherwise
the required correction is obtained by rounding each necessary coefficient
ratio upward to a multiple of $2^{-32}$ and taking the largest correction.
On the interior of the cell this proves
\[
 \frac{N_\pi}{z_\pi^k}
 =J_\pi-c_\pi+
   \frac{O_\pi+c_\pi Q_\pi}{z_\pi^k}
 \ge J_\pi-c_\pi.
\]
The cell is accepted when the homogeneous coefficients of
$K+\sum_\pi(J_\pi-c_\pi)$ are all nonnegative, after degree elevation.
The polynomials $J_\pi$ in the source are obtained by rational interpolation
on the degree-three simplex lattice.  Their interpolation accuracy is not
assumed: the full remainder is certified by the displayed identity.

Unaccepted cells are bisected at an edge midpoint.  If vertex $i$ is
replaced by the midpoint of vertices $i,j$, the barycentric substitution is
\[
 X_i=X_i'/2,\qquad X_j=X_j'+X_i'/2.
\]
The binomial expansion gives exact rational coefficients for the child
cell.  The two choices of the replaced endpoint cover the original
simplex, corresponding to $X_i\le X_j$ and $X_j\le X_i$.  Repeating this
construction therefore proves coverage, rather than merely testing a set
of posterior points.

The full integer verifier of \cite{OpenAI4}, with assertions enabled, was
replayed for both targets.  It terminated with the following counts and no
unresolved cells:
\begin{center}
\begin{tabular}{crrrr}
\hline
Target & Visited cells & Maximum depth & Taylor leaves & Remainder leaves\\
\hline
$a$ & 4,479 & 19 & 166 & 2,074\\
$b$ & 1,935 & 20 & 97 & 871\\
\hline
\end{tabular}
\end{center}
The scalar coefficients and the twenty-four induced coordinate actions were
also checked independently by rational arithmetic.  All sign decisions in
the certificates use integer numerators and positive exact denominators.

Finally, denominator clearing and the lower bounds above are used first on
cell interiors, where every $z_\pi>0$.  The weighted terms
$z_\pi A(P^\pi)$ and $z_\pi H_4(P^\pi)$ converge to zero when
$z_\pi\to0$, since $A,H_4$ are bounded on $\Delta_4$.  They therefore have
continuous extensions to every boundary point.  This proves the boundary
cases without dividing by a vanishing denominator.  The original source
contains the complete verifier and its deterministic subdivision rule.

\subsection{Closure of the four-state cone}

Let $\mathcal C_4$ consist of the permutation-symmetric posterior laws
$\mu$ satisfying $\mathbb E_\mu H_4\ge0$.  The condition is imposed on the
law, not on each posterior in its support.  Full revelation belongs to the
cone because $H_4(e_i)=13/100$; the uninformative experiment belongs because
$H_4(u_4)=0$.

\begin{lemma}[Four-state experiment closure]
\label{lem:q4-cone-closure}
The class $\mathcal C_4$ is preserved by a channel edge with
$0\le\lambda\le1$, by conditionally independent products, and by observed
mixtures with an input-independent mixing label.  If
$\mu,\nu\in\mathcal C_4$, and
$a=\mathbb E_\mu A$, $b=\mathbb E_\nu A$, then
\begin{equation}
 \mathbb E_{\mu\star\nu}A
 \le a+b-
 \frac{(\mathbb E_\mu A^2)b+a(\mathbb E_\nu A^2)}{1000}
 \le a+b-\frac{ab(a+b)}{1000}.
 \label{eq:q4-experiment-saving}
\end{equation}
\end{lemma}

\begin{proof}
An edge sends $p$ to $u_4+\lambda(p-u_4)$, and hence sends $Lp$ to
$\lambda Lp$.  Equation~\eqref{eq:q4-radial} preserves the expectation
condition, and $A$ is multiplied by $\lambda^2$.

For the product of two experiments, the expectation of a posterior
functional $f$ is
\[
 \mathbb E_{\mu\star\nu}f
 =\mathbb E_{\mu\otimes\nu}[z f(P)].
\]
The laws in the reference measure $\mu\otimes\nu$ are independent; the
factor $z$ gives the actual joint observation law.  By permutation
symmetry, averaging the integrand over the relative permutations does not
change the expectation.  Integrating \eqref{eq:q4-pair-H} therefore gives
\[
 \mathbb E_{\mu\star\nu}H_4
 \ge (\mathbb E_\mu H_4)(\mathbb E_\nu G_4)
      +(\mathbb E_\mu G_4)(\mathbb E_\nu H_4)\ge0.
\]
The product is symmetric as well.  Integrating
\eqref{eq:q4-pair-A} and discarding the nonnegative products involving
$F_4$ gives the first inequality in
\eqref{eq:q4-experiment-saving}.  Jensen's inequality gives the second.
Observed input-independent mixtures preserve both symmetry and the
expectation condition by linearity.
\end{proof}

\section{Nonreconstruction for arbitrary finite-mean offspring}
\label{sec:finite-mean-nonreconstruction}

Let $D$ be a nonnegative, almost surely finite integer-valued random
variable with $d=\mathbb E D<\infty$, and let $T$ be the associated
Galton--Watson tree.  The tree and all auxiliary child orderings are
independent of the spins.  The observer knows the tree and the spins at
generation $n$.  Knowing the continuation of the tree beyond that
generation adds no spin information conditional on the observed finite
prefix.

For $q=3,4$, write $p_n$ for the posterior of the uniformly distributed
root spin.  Use the quadratic statistics
\[
 J_3(p)=\frac{3\sum_{i=1}^3p_i^2-1}{2},\qquad
 J_4(p)=4\sum_{i=1}^4p_i^2-1,
 \qquad m_n^{(q)}=\mathbb E J_q(p_n).
\]
Thus $J_3=x$ in the three-state section and $J_4=A$ in the preceding
section.  Initially $m_0^{(3)}=1$ and $m_0^{(4)}=3$.  Set
\begin{equation}
 \rho=d\lambda^2,\qquad
 \delta=\mathbb E(D-1)_+=d-\mathbb P(D\ge1).
 \label{eq:gw-rho-delta}
\end{equation}
Only the first offspring moment is involved in $\delta$.

\begin{proposition}[Quadratic and cubic drift]
\label{prop:gw-drift}
For $0\le\lambda\le1$ and every $n\ge0$,
\begin{align}
 m_{n+1}^{(3)}
 &\le \rho m_n^{(3)}-c_3\bigl(m_n^{(3)}\bigr)^2,
 &c_3&=\frac{\delta\lambda^4}{1000},
 \label{eq:gw-q3-drift}\\
 m_{n+1}^{(4)}
 &\le \rho m_n^{(4)}-c_4\bigl(m_n^{(4)}\bigr)^3,
 &c_4&=\frac{\delta\lambda^6}{500}.
 \label{eq:gw-q4-drift}
\end{align}
\end{proposition}

\begin{proof}
At depth zero the fully revealing experiment belongs to the appropriate
cone.  Edge attenuation, conditionally independent products, and observed
offspring mixtures show inductively that every finite-depth experiment
belongs to that cone.  These closures also apply to every partial product
of incoming branch experiments.

Condition on $D=k$ and put $b=\lambda^2m_n^{(q)}$.  Every child branch
has mean statistic $b$, since one channel edge multiplies the quadratic
statistic by $\lambda^2$.  If $j\ge1$ branches have already been observed,
their partial product has mean statistic at least $b$.  To see this,
let $P^{(j)}$ be the posterior after those $j$ branches and $P^{(1)}$ the
posterior after the first branch.  The tower property gives
\[
 P^{(1)}=\mathbb E[P^{(j)}\mid\hbox{first branch observation}],
\]
and conditional Jensen applies because both $J_3$ and $J_4$ are convex.

The three-state pair inequality proved earlier gives, for a partial
product with mean $a\ge b$ and one further branch,
\[
 m(\hbox{new product})\le a+b-\frac{ab}{1000}
 \le a+b-\frac{b^2}{1000}.
\]
Applying this at each of the $k-1$ fusions, and including $k=0,1$, yields
\[
 \mathbb E[J_3(p_{n+1})\mid D=k]
 \le kb-\frac{(k-1)_+b^2}{1000}.
\]
For four states, \eqref{eq:q4-experiment-saving} instead gives a saving
of at least
\[
 \frac{ab(a+b)}{1000}\ge\frac{b^3}{500}
\]
at each fusion.  Therefore
\[
 \mathbb E[J_4(p_{n+1})\mid D=k]
 \le kb-\frac{(k-1)_+b^3}{500}.
\]
Averaging over the observed offspring count proves both recursions.
The expectation is finite because $\mathbb E D<\infty$.
\end{proof}

\begin{proposition}[Decay at and below the Kesten--Stigum value]
\label{prop:gw-decay}
Assume $0<\lambda<1$ and $\rho\le1$.  Let
\[
 a_n^{(q)}=\mathbb E\operatorname{TV}(p_n,u_q).
\]
If $\rho<1$, then
\begin{align}
 m_n^{(3)}&\le\rho^n,
 &a_n^{(3)}&\le\frac{\rho^{n/2}}{\sqrt2},
 \label{eq:q3-subcritical-rate}\\
 m_n^{(4)}&\le3\rho^n,
 &a_n^{(4)}&\le\frac{\sqrt3}{2}\rho^{n/2}.
 \label{eq:q4-subcritical-rate}
\end{align}
If $\rho=1$, then $c_3,c_4>0$ and
\begin{align}
 m_n^{(3)}&\le\frac1{1+c_3n},
 &a_n^{(3)}&\le\frac1{\sqrt{2(1+c_3n)}},
 \label{eq:q3-critical-rate}\\
 m_n^{(4)}&\le(1/9+2c_4n)^{-1/2},
 &a_n^{(4)}&\le\frac12(1/9+2c_4n)^{-1/4}.
 \label{eq:q4-critical-rate}
\end{align}
In both cases nonreconstruction holds for almost every realized tree.
\end{proposition}

\begin{proof}
Below criticality, omit the negative terms in
\eqref{eq:gw-q3-drift}--\eqref{eq:gw-q4-drift} and iterate.  The
total-variation estimates follow from Cauchy--Schwarz:
\[
 \operatorname{TV}(p,u_3)\le\sqrt{J_3(p)/2},\qquad
 \operatorname{TV}(p,u_4)\le\frac12\sqrt{J_4(p)},
\]
followed by Jensen's inequality.

At $\rho=1$, the condition $\lambda<1$ implies $d>1$, and hence
$\delta=d-\mathbb P(D\ge1)>0$.  For three states, while
$m_n=m_n^{(3)}>0$,
\[
 \frac1{m_{n+1}}
 \ge\frac1{m_n(1-c_3m_n)}
 \ge\frac1{m_n}+c_3.
\]
The denominator is positive: $m_n\le1$ and
$c_3\le d\lambda^4/1000\le1/1000$ when $\rho\le1$.
If an iterate is zero, all subsequent iterates are zero and the stated
bound is immediate.  Otherwise iteration from $m_0=1$ proves
\eqref{eq:q3-critical-rate}.

For four states, write $m_n=m_n^{(4)}$.  Since $m_n\le3$ and
$c_4m_n^2\le9d\lambda^6/500\le9/500<1$, the analogous calculation gives
\[
 \frac1{m_{n+1}^2}
 \ge\frac1{m_n^2}(1-c_4m_n^2)^{-2}
 \ge\frac1{m_n^2}+2c_4.
\]
Iteration from $m_0=3$, with the same zero-iterate convention, proves
\eqref{eq:q4-critical-rate}.

For a fixed locally finite tree $T$, put
\[
 a_n^{(q)}(T)=
 \mathbb E[\operatorname{TV}(p_n,u_q)\mid T].
\]
The root spin, level $n$, and level $n+1$ form a Markov chain, so
$a_n^{(q)}(T)$ is nonincreasing by data processing.  Its nonnegative limit
therefore exists.  The annealed estimates show that its expectation tends
to zero.  Bounded convergence yields
\[
 \mathbb{E}\bigl[\lim_{n\to\infty}a_n^{(q)}(T)\bigr]=0,
\]
so the limit is zero for almost every tree.  The conclusion also holds
almost surely conditional on survival whenever survival has positive
probability.  No independence assumption after conditioning on survival
is needed.
\end{proof}

\section{Reconstruction above the threshold}
\label{sec:finite-mean-supercritical}

The argument in this section works for every finite number of states
$q\ge2$.  It supplies the supercritical direction with only a finite first
offspring moment.

\begin{proposition}[Bounded pruning]
\label{prop:bounded-pruning}
Let $\mathbb E D=d<\infty$ and $0<\lambda<1$.  If
$d\lambda^2>1$, then the annealed reconstruction advantage is bounded
away from zero uniformly in depth.
\end{proposition}

\begin{proof}
By monotone convergence, choose $L<\infty$ such that
\[
 D_L=\min(D,L),\qquad \mu=\mathbb E D_L,
 \qquad \mu\lambda^2>1.
\]
Keep the first $D_L$ children at every retained vertex, using a child
ordering independent of the spins.  The retained subtree has iid bounded
offspring $D_L$, and its depth-$n$ spins are part of the original
depth-$n$ observation.

Choose a channel eigenfunction $\phi$ with uniform mean zero and second
moment one; for example,
\[
 \phi(1)=\sqrt{q-1},\qquad
 \phi(i)=-\frac1{\sqrt{q-1}}\quad(2\le i\le q).
\]
It satisfies
$\mathbb E[\phi(\sigma_{\rm child})\mid\sigma_{\rm parent}]
=\lambda\phi(\sigma_{\rm parent})$ and
$|\phi|\le B_q:=\sqrt{q-1}$.
On the retained tree define
\[
 Z_n=(\mu\lambda)^{-n}
       \sum_{|v|=n}\phi(\sigma_v).
\]
The first-moment recursion gives
\begin{equation}
 \mathbb E[Z_n\mid\sigma_{\rm root}]
 =\phi(\sigma_{\rm root}),\qquad
 \mathbb E[\phi(\sigma_{\rm root})Z_n]=1.
 \label{eq:pruning-first-moment}
\end{equation}

Let $M_n=\mathbb E Z_n^2$ and
$\alpha_L=\mathbb E[D_L(D_L-1)]/\mu^2$.  Writing the depth-$(n+1)$
sum as a sum of child contributions gives
\begin{equation}
 M_{n+1}=\frac{M_n}{\mu\lambda^2}+\alpha_L,
 \qquad M_0=1.
 \label{eq:pruning-second-moment}
\end{equation}
Indeed the diagonal terms have total expectation $\mu M_n$.  For distinct
children, conditional independence given the root and
\eqref{eq:pruning-first-moment} give product expectation
$\lambda^2\phi(\sigma_{\rm root})^2$, whose mean is $\lambda^2$.
Dividing by $(\mu\lambda)^2$ proves
\eqref{eq:pruning-second-moment}.  Since $D_L$ is bounded and
$r_L=(\mu\lambda^2)^{-1}<1$,
\[
 M_n\le C_L:=1+\frac{\alpha_L}{1-r_L}<\infty
 \qquad\hbox{for all }n.
\]

Let $\mathcal F_n$ contain the original tree, its level-$n$ spins, and
the auxiliary ordering.  The statistic $Z_n$ is $\mathcal F_n$-measurable.
Put $\eta_n=\mathbb E[\phi(\sigma_{\rm root})\mid\mathcal F_n]$.
Cauchy--Schwarz and \eqref{eq:pruning-first-moment} imply
\[
 1=\mathbb E[\eta_nZ_n]
 \le\bigl(\mathbb E\eta_n^2\,C_L\bigr)^{1/2},
 \qquad \mathbb E\eta_n^2\ge C_L^{-1}.
\]
The ordering is independent of the root and spins conditional on the
original observed tree.  Consequently adjoining it does not change the
root posterior.  For that posterior,
\[
 |\eta_n|\le B_q,\qquad
 \operatorname{TV}(p_n,u_q)\ge\frac{|\eta_n|}{2B_q}.
\]
Therefore
\begin{equation}
 a_n^{(q)}
 \ge\frac{\mathbb E|\eta_n|}{2B_q}
 \ge\frac{\mathbb E\eta_n^2}{2B_q^2}
 \ge\frac1{2(q-1)C_L}>0.
 \label{eq:supercritical-advantage}
\end{equation}
Only the truncated offspring variable was required to have a second
moment.
\end{proof}

\begin{lemma}[The reconstruction event on a Galton--Watson tree]
\label{lem:survival-reconstruction}
For $0<\lambda<1$ and $d\lambda^2>1$, reconstruction holds for almost
every surviving Galton--Watson tree and fails on the extinction event.
\end{lemma}

\begin{proof}
First fix a deterministic locally finite rooted tree.  Its root
reconstructs if and only if at least one of its child subtrees reconstructs.
If a child reconstructs, observing only that child's descendants gives
the root posterior
\[
 u_q+\lambda(p_{\rm child}-u_q).
\]
Its total-variation advantage is $\lambda$ times the child's advantage,
so remains positive in the limit.

Conversely, suppose all child subtrees are nonreconstructing.  For a finite
input alphabet, vanishing expected posterior total variation implies
vanishing mutual information: the posterior relative entropy from the
uniform prior is bounded and continuous on the compact simplex.  Thus
every child's boundary observation has vanishing mutual information with
the child spin and, by data processing, with the root spin.  There are
only finitely many children.  Their observations $Y_1,\ldots,Y_k$ are
independent conditional on the root, and hence
\[
 I(\sigma_{\rm root};Y_1,\ldots,Y_k)
 \le\sum_{j=1}^k I(\sigma_{\rm root};Y_j)\longrightarrow0.
\]
Here the inequality follows from subadditivity of unconditional entropy
and additivity of conditional entropy.  Pinsker's inequality then implies
nonreconstruction at the root.

Let $r_* $ be the probability that a Galton--Watson tree reconstructs.
The preceding equivalence and independence of child subtrees yield
\[
 r_*=1-\mathbb E(1-r_*)^D.
\]
Let $\xi$ be the extinction probability.  Since
$d\lambda^2>1$ implies $d>1$, the offspring generating function has
exactly the two fixed points $\xi$ and $1$ in $[0,1]$.  For completeness,
it is strictly convex on $(0,1)$ because $\mathbb P(D\ge2)>0$;
its derivative at $1$ equals $d>1$, and $\xi$ is its least fixed point.
These facts exclude any third fixed point.

The positive lower bound \eqref{eq:supercritical-advantage}, together
with bounded convergence for the fixed-tree advantage, implies $r_*>0$.
Thus $1-r_*=\xi$.  Reconstruction is impossible on an extinct tree,
whose boundary is eventually empty.  Consequently the reconstruction
event agrees with survival up to a null set.
\end{proof}

\section{The finite-mean dichotomy and endpoint conventions}
\label{sec:full-gw-dichotomy}

\begin{theorem}[Three and four states on finite-mean Galton--Watson trees]
\label{thm:full-finite-mean-dichotomy}
Let $D$ be a nonnegative, almost surely finite integer-valued random
variable with finite mean $d$.  Consider the symmetric ferromagnetic
$q$-state broadcast process on its observed Galton--Watson tree, with
$q\in\{3,4\}$ and $0<\lambda<1$.
If $d\lambda^2\le1$, nonreconstruction holds for almost every tree.
If $d\lambda^2>1$, reconstruction holds for almost every surviving
tree and fails on extinct trees.  In particular, annealed reconstruction
occurs if and only if
\[
 d\lambda^2>1.
\]
No offspring moment beyond $\mathbb E D<\infty$ is assumed.  At equality,
the quantitative bounds are \eqref{eq:q3-critical-rate} and
\eqref{eq:q4-critical-rate}.
\end{theorem}

\begin{proof}
Proposition~\ref{prop:gw-decay} proves the assertion at and below equality.
Proposition~\ref{prop:bounded-pruning} and
Lemma~\ref{lem:survival-reconstruction} prove the assertions above equality.
\end{proof}

\begin{remark}[Regular trees and endpoints]
Taking $D\equiv d$ gives every regular tree of integer degree $d\ge2$.
At $\lambda=0$, all positive-depth observations are independent of the
root.  At $\lambda=1$, every descendant copies the root, so reconstruction
is perfect on survival and impossible after extinction.  A subcritical
or nondegenerate critical finite-mean Galton--Watson process dies out
almost surely.  The exceptional law $D\equiv1$ instead gives an infinite
noiseless line: at $\lambda=1$ it reconstructs perfectly although
$d\lambda^2=1$.  A statement including all closed endpoints must record
this exception.  The argument concerns nonnegative channel eigenvalues;
it makes no assertion about the antiferromagnetic case.
\end{remark}


\section*{Acknowledgments and funding}
The research presented in this paper was supported by the European Research Council (ERC) under the European Union's Horizon Europe research and innovation programme (grant agreement No.~101041711), by the Simons Foundation as part of the Collaboration on the Mathematical and Scientific Foundations of Deep Learning, by Heights Labs, by Convex Nexus Capital, by the Israel Science Foundation (grant number 2258/19), and by the Israel Science Foundation (ISF Grant 4101/25).

\section*{Declaration of competing interest}
The author holds equity in Heights Labs and in Convex Nexus Capital and has a management role at Convex Nexus Capital. These entities are included in the funding declaration above.

\section*{Declaration of generative AI and AI-assisted technologies}
The three-state argument, the finite-mean offspring deductions, the literature search, the verification code and the manuscript were produced by OpenAI AI systems under the author's direction. The four-state posterior-cone theorem is prior work and is attributed explicitly in the text. The computations were checked by exact certificate reruns and an independent reconstruction of all three-state polynomial coefficients. The author selected the problem and commissioned and directed the work; the author did not carry out the mathematics or the writing. Responsibility for any decision to submit or publish the manuscript rests with the author.


\appendix
% Reproducibility appendix. Requires amsmath; citation key: OpenAI4.
\section{Exact coefficient certificates and reproducibility}
\label{app:exact-certificates}

This appendix specifies the finite arithmetic certificates used in the proof.
The three-state inequalities are verified by two different polynomial
constructions.  The four-state inequalities use the published certificate
of \cite{OpenAI4}, whose two complete traversals were rerun.
All sign decisions use integer or rational arithmetic.

\subsection{The three-state residuals and their complete domain}

For \(p\in\Delta_3\), retain the definitions
\[
 x(p)=\frac16\sum_i(3p_i-1)^2,\qquad
 y(p)=\frac16\sum_i(3p_i-1)^3,\qquad
 H=y-\frac{x^2}{2},\quad F=\frac32x-y,\quad G=1-x.
\]
For \(\pi\in S_3\), put
\[
 z_\pi=3\sum_i p_iq_{\pi(i)},\qquad
 r_{\pi,i}=\frac{3p_iq_{\pi(i)}}{z_\pi},\qquad
 \epsilon=\frac1{1000}.
\]
A term \(z_\pi f(r_\pi)\) is defined as zero when \(z_\pi=0\).
The two residuals certified below are
\begin{align}
 R_A(p,q)
 &=x(p)+x(q)-\frac16\sum_{\pi\in S_3}z_\pi x(r_\pi)\notag\\
 &\quad-H(p)F(q)-F(p)H(q)-\epsilon x(p)x(q), \label{eq:app-RA}\\
 R_B(p,q)
 &=\frac16\sum_{\pi\in S_3}z_\pi H(r_\pi)\notag\\
 &\quad-H(p)G(q)-G(p)H(q)-\epsilon x(p)x(q). \label{eq:app-RB}
\end{align}

It suffices to work in the ordered chambers for \(p\) and \(q\):
all scalar functions are symmetric, and the fusion terms average over
every relative permutation.  Introduce nonnegative variables
\((a,b,c)\), \((d,e,f)\), and let \(S=a+b+c>0\), \(T=d+e+f>0\).  Set
\[
 p=\frac{(6a+3b+2c,\;3b+2c,\;2c)}{6S},
 \qquad
 q=\frac{(6d+3e+2f,\;3e+2f,\;2f)}{6T}.
\]
This covers the whole ordered simplex.  Indeed, for
\(p_1\ge p_2\ge p_3\), the normalized coordinates
\[
 (a,b,c)=(p_1-p_2,\;2(p_2-p_3),\;3p_3)
\]
are nonnegative and sum to one.  Thus the certificate is not restricted
to sampled posterior vectors or to a neighborhood of the uniform law.

\subsection{An explicit integer polynomial certificate}

Define
\[
 X=4a^2+2ab+b^2,\qquad
 Y=8a^3+6a^2b-3ab^2-b^3,
\]
and use primes for the same expressions in \((d,e,f)\).  Direct
substitution gives
\[
 x(p)=\frac{X}{4S^2},\qquad y(p)=\frac{Y}{8S^3}.
\]
Writing
\[
 \widehat H=4YS-X^2,\qquad
 \widehat F=3XS-Y,\qquad
 \widehat G=4S^2-X,
\]
we have
\[
 H(p)=\frac{\widehat H}{32S^4},\qquad
 F(p)=\frac{\widehat F}{8S^3},\qquad
 G(p)=\frac{\widehat G}{4S^2}.
\]
The auxiliary signs used in the expectation argument follow, for example,
from the explicit expansions
\begin{align*}
 \widehat F
 &=4a^3+12a^2b+12ab^2+4b^3+12a^2c+6abc+3b^2c,\\
 \widehat G
 &=6ab+8ac+3b^2+8bc+4c^2.
\end{align*}

Let \(P=(6a+3b+2c,3b+2c,2c)\) and
\(Q=(6d+3e+2f,3e+2f,2f)\).  For each permutation define
\[
 U_{\pi,i}=P_iQ_{\pi(i)},\quad
 Z_\pi=\sum_iU_{\pi,i},\quad
 W_{\pi,i}=3U_{\pi,i}-Z_\pi,
\]
\[
 V_\pi=\sum_iW_{\pi,i}^{\,2},\qquad
 J_\pi=36Z_\pi W_{\pi,1}W_{\pi,2}W_{\pi,3}-V_\pi^2.
\]
Since \(\sum_iW_{\pi,i}=0\), the weighted fusion terms satisfy
\[
 z_\pi=\frac{Z_\pi}{12ST},\qquad
 x(r_\pi)=\frac{V_\pi}{6Z_\pi^2},\qquad
 H(r_\pi)=\frac{J_\pi}{72Z_\pi^4}.
\]
In particular, the factor \(z_\pi\) from the posterior product law reduces
the denominator powers to one and three.

Set \(D_A=\prod_\pi Z_\pi\), \(D_B=\prod_\pi Z_\pi^3\), and define
\begin{align*}
 K_A={}&8000\bigl(XS^2T^4+X'S^4T^2\bigr)
       -125\bigl(\widehat H\widehat F'T+\widehat F\widehat H'S\bigr)
       -2XX'S^2T^2,\\
 K_B={}&-125\bigl(\widehat H\widehat G'T^2+
                         \widehat G\widehat H'S^2\bigr)-XX'S^2T^2.
\end{align*}
The following two polynomials have integer coefficients:
\begin{align}
 \mathcal P_A
 &=27K_AD_A-2000S^3T^3
       \sum_\pi V_\pi\prod_{\rho\ne\pi}Z_\rho, \label{eq:app-PA}\\
 \mathcal P_B
 &=81K_BD_B+250S^3T^3
       \sum_\pi J_\pi\prod_{\rho\ne\pi}Z_\rho^3. \label{eq:app-PB}
\end{align}
Their exact relation to the desired residuals is
\begin{equation}
 \mathcal P_A=864000S^4T^4D_A R_A,\qquad
 \mathcal P_B=1296000S^4T^4D_B R_B.                 \label{eq:app-clearing}
\end{equation}

Expanding \eqref{eq:app-PA}--\eqref{eq:app-PB} in the six nonnegative
barycentric variables gives the complete coefficient inventory:
\begin{center}
\begin{tabular}{lrrrrr}
\hline
Certificate & Bidegree & Positions & Positive & Zero & Negative\\
\hline
A & \((10,10)\) & \(4\,356\) & \(3\,873\) & \(483\) & \(0\)\\
B & \((22,22)\) & \(76\,176\) & \(71\,627\) & \(4\,549\) & \(0\)\\
\hline
\end{tabular}
\end{center}
The number of positions is \(\binom{n+2}{2}^2\) in bidegree \((n,n)\).
Every coefficient is checked; no subdivision is needed for these
three-state polynomials.

Nonnegative coefficients imply \(\mathcal P_A,\mathcal P_B\ge0\).
For interior probability vectors, \(Z_\pi>0\) for all \(\pi\), so
\eqref{eq:app-clearing} proves \(R_A,R_B\ge0\).
Boundary points are included by continuity: \(x\) and \(H\) are bounded on
the probability simplex, hence \(z_\pi x(r_\pi)\) and
\(z_\pi H(r_\pi)\) tend to zero when \(z_\pi\) tends to zero.
Approximation by interior pairs therefore establishes
\eqref{eq:app-RA}--\eqref{eq:app-RB} on the entire domain.

\subsection{Independent verification}

The source \path{analytic/q3_exact_pair_simple.cpp} constructs rational
homogeneous polynomials directly from the definitions of the posterior
invariants.  Separately,
\path{computation/q3_independent_coefficients.cpp} expands
\eqref{eq:app-PA}--\eqref{eq:app-PB} in a sparse six-variable integer
representation.  Its construction does not use the first encoder's
invariants, homogenization, tensor multiplication, or rational-polynomial
operations.  Only after constructing its polynomials does it call the
first encoder to compare the complete coefficient arrays.

If \(\mathcal C_A,\mathcal C_B\) denote the rational homogeneous outputs of
the first encoder, every coefficient satisfies the exact identities
\[
 \mathcal P_A=144000\cdot12^6\,\mathcal C_A,\qquad
 \mathcal P_B=216000\cdot12^{18}\,\mathcal C_B.
\]
The comparison uses integer cross-products with the stored denominators.
Both constructions use arbitrary-precision integers.  The independent
implementation packs six exponents into five-bit fields and checks that
each group's total degree is below \(32\); the maximum used is \(22\),
so exponent addition cannot carry between fields.

Both implementations completed with assertions enabled and
UndefinedBehaviorSanitizer.  The clean certificates accept at their root,
so floating-point edge-selection coordinates in the first encoder are
never consulted.  No sampled inequality or floating-point sign decision
is used.

The supplementary files contain both implementations, the complete coefficient tables, and the verification output. The independent program regenerates every coefficient, including all zeros.

\subsection{Four-state certificate}

The four-state pair inequalities and their certificate algorithm are those of \cite{OpenAI4}; the supplementary files include the unchanged verifier from the version linked in that reference.

The verifier works on products of ordered posterior simplices and sums
over all \(24\) relative colour permutations.  It certifies subdomains
using either a Taylor lower bound for reciprocal powers with a
nonnegative multiplier, or cubic interpolation with a rigorously bounded
cleared polynomial remainder.  Exact midpoint bisection covers a
subdomain whenever neither certificate yet applies.  A successful
traversal certifies every resulting subdomain; see \cite{OpenAI4} for the
inequalities and the justification of these two tests.

Both full traversals were executed with the published C++ source unchanged,
exact integer/rational arithmetic, and assertions enabled:
\begin{center}
\begin{tabular}{lrrrr}
\hline
Mode & Nodes & Max.\ depth & Interpolation leaves & Taylor leaves\\
\hline
A & \(4\,479\) & \(19\) & \(2\,074\) & \(166\)\\
B & \(1\,935\) & \(20\) & \(871\) & \(97\)\\
\hline
\end{tabular}
\end{center}
Both exited with status zero.  Their terminal lines were
\begin{verbatim}
done 4479 19 lp=2074 lt=166
done 1935 20 lp=871 lt=97
\end{verbatim}
Complete output records are included with the supplementary code.

Additional checks verify the four homogeneous scalar polynomials' coefficient minima, all \(24\) symmetry actions, and the exact arithmetic, interpolation, and subdivision operations under UndefinedBehaviorSanitizer.  The scalar coefficient minima are
\(19/50,0,87/1000,0\) for \(H+1/2,F,G,(3-\mathrm{Euler})H\),
respectively; these are coefficient minima, not asserted minima of the
functions.  The lightweight checks are supplemented by the two complete traversals reported above.

\subsection{Compact reproduction instructions}

The accompanying source archive preserves the cited files and their
relative paths.  Requirements are Python~3, a C++17 compiler with the
GCC-compatible headers used by the published verifier, and Boost
multiprecision headers.  Keep assertions enabled.  If Boost is outside
the compiler's default include path, set \texttt{CPLUS\_INCLUDE\_PATH}
to its include directory.  From the archive root:
\begin{verbatim}
cd computation
g++ -std=c++17 -O2 -UNDEBUG -fsanitize=undefined \
  -fno-sanitize-recover=all \
  ../analytic/q3_exact_pair_simple.cpp -o q3_check
./q3_check 0
./q3_check 1
g++ -std=c++17 -O2 -UNDEBUG -fsanitize=undefined \
  -fno-sanitize-recover=all \
  q3_independent_coefficients.cpp -o q3_independent
./q3_independent 0
./q3_independent 1
cd published_q4/verification
python3 checks/verify.py
g++ -std=c++17 -O3 -UNDEBUG \
  checks/certificate.cpp -o artifacts/certificate
./artifacts/certificate
./artifacts/certificate b
\end{verbatim}
Every command must exit with status zero.  The q3 outputs must reproduce
the coefficient counts above, and the q4 executions must print the two
stated completion lines. The probabilistic deductions from these algebraic inequalities are supplied in the main
text.



\section*{Code availability}

The supplementary material contains the complete \LaTeX{} source, both three-state verification programs, the full coefficient tables, and the four-state verifier of \cite{OpenAI4}, with instructions for reproducing the exact checks. The version of the four-state source is identified in the bibliography.

\begin{thebibliography}{9}
\bibitem{Sly2011}
A.~Sly, \emph{Reconstruction for the Potts model}, Annals of Probability \textbf{39} (2011), 1365--1406. Earlier preprint: \emph{Reconstruction of symmetric Potts models}, \href{https://arxiv.org/abs/0811.1208}{arXiv:0811.1208}.

\bibitem{OpenAI3}
OpenAI, \emph{The exact reconstruction threshold for the three-state symmetric channel}, manuscript dated 25 September 2026.
\href{https://github.com/openai/math/tree/1581d4cbdd69a0faf34b2089614fe2a814316a6a/preprints/The-exact-reconstruction-threshold-for-the-three-state-symmetric-channel-September-25-2026}{Manuscript and source files}.

\bibitem{OpenAI4}
OpenAI, \emph{The Reconstruction Threshold for the Ferromagnetic Four-State Potts Model}, manuscript dated 5 October 2026.
\href{https://github.com/openai/math/tree/1581d4cbdd69a0faf34b2089614fe2a814316a6a/preprints/The-Reconstruction-Threshold-for-the-Ferromagnetic-Four-State-Potts-Model-October-5-2026}{Manuscript and exact-certificate source}.
\end{thebibliography}
\end{document}
